\documentclass[1p,times]{elsarticle}
\usepackage{amsmath,amssymb}
\usepackage{booktabs}
\usepackage{graphicx}
\usepackage[colorlinks=true,linkcolor=blue,citecolor=blue,urlcolor=blue,hypertexnames=false]{hyperref}
\usepackage{microtype}
\usepackage{placeins}

\providecommand{\orcid}[1]{\href{https://orcid.org/#1}{\texttt{#1}}}

\journal{Applied Soft Computing}

\begin{document}

\begin{frontmatter}
\title{Supplementary Material: Correctness and Predictive Uncertainty in Bayesian TSK Fuzzy Regression}
\author[addr]{Zhiren Xiao~\orcid{0009-0008-2164-4557}}
\address[addr]{Guangdong University of Finance, Guangzhou, China}
\end{frontmatter}

\section{Rebuilt Protocol and Provenance}

All tables in this supplement are generated from the JSON outputs in the
date-stamped rebuild-results directory.  The main experiment uses 30
outer 80/20 train--test splits with split seeds starting at 42, five fuzzy
rules, fuzzy $c$-means exponent $m=2$, standardized features, and Gaussian
membership spreads scaled by $s=1.5$.  The rule-level Gibbs sampler uses
$1{,}000$ burn-in iterations and $2{,}000$ retained draws; coefficient-level
SSVS uses $800$ burn-in and $800$ retained draws.  Bayesian intervals are
the direct 2.5th and 97.5th posterior-predictive quantiles.  The conformal
baseline uses an inner calibration split and is therefore reported
separately.

The reported protocol records every split-level metric, runtime, active-rule
count, and available effective sample size.  The generated tables and
figures are local working artifacts; no new repository release, DOI, or
external submission was created in this work.

\section{Mean Absolute Error}

\begin{table}[htbp]
\centering
\caption{Mean absolute error under the rebuilt 30-split protocol.}
\label{tab:mae}
\scriptsize
\setlength{\tabcolsep}{3pt}
\begin{tabular}{lrrr}
\toprule
Method & Energy-Heating & Energy-Cooling & Concrete \\
\midrule
TSK-LS & 1.521 & 1.926 & 5.828 \\
Bayesian-TSK & 1.526 & 1.931 & 5.829 \\
TSK-SpikeSlab-BIC & 1.539 & 1.938 & 5.829 \\
TSK-SpikeSlab-Gibbs & 1.531 & 1.941 & 5.849 \\
Random Forest & 0.329 & 1.014 & 3.617 \\
SVR & 1.803 & 2.054 & 7.767 \\
Gaussian Process & 0.338 & 0.924 & 4.009 \\
Conformal-TSK-LS & 1.563 & 1.957 & 5.793 \\
\bottomrule
\end{tabular}
\end{table}


\FloatBarrier

\section{Posterior-Predictive Metrics}

\begin{table}[htbp]
\centering
\caption{Posterior-predictive metrics for methods that produce probabilistic intervals. WIS and CRPS are unavailable for the conformal baseline because it does not define a posterior-predictive distribution.}
\label{tab:predictive}
\scriptsize
\setlength{\tabcolsep}{3pt}
\begin{tabular}{llrrrr}
\toprule
Target & Method & PICP & MPIW & WIS & CRPS \\
\midrule
Energy-Heating & Bayesian-TSK & 0.940 & 8.164 & 0.938 & 1.129 \\
Energy-Heating & TSK-SpikeSlab-BIC & 0.940 & 8.151 & 0.944 & 1.138 \\
Energy-Heating & TSK-SpikeSlab-Gibbs & 0.940 & 8.180 & 0.939 & 1.132 \\
Energy-Heating & Gaussian Process & 0.931 & 1.748 & 0.210 & 0.252 \\
Energy-Cooling & Bayesian-TSK & 0.932 & 10.270 & 1.203 & 1.442 \\
Energy-Cooling & TSK-SpikeSlab-BIC & 0.934 & 10.248 & 1.205 & 1.445 \\
Energy-Cooling & TSK-SpikeSlab-Gibbs & 0.934 & 10.295 & 1.203 & 1.443 \\
Energy-Cooling & Gaussian Process & 0.938 & 5.175 & 0.597 & 0.706 \\
Concrete & Bayesian-TSK & 0.938 & 28.943 & 3.505 & 4.248 \\
Concrete & TSK-SpikeSlab-BIC & 0.938 & 28.894 & 3.505 & 4.249 \\
Concrete & TSK-SpikeSlab-Gibbs & 0.938 & 29.148 & 3.518 & 4.261 \\
Concrete & Gaussian Process & 0.939 & 20.457 & 2.447 & 2.946 \\
\bottomrule
\end{tabular}
\end{table}


\begin{table}[htbp]
\centering
\caption{Split-conformal TSK-LS results. Calibration is performed inside each outer training split, so this row is not a posterior-predictive baseline under the same protocol as Table~\ref{tab:predictive}.}
\label{tab:conformal}
\scriptsize
\begin{tabular}{lrrr}
\toprule
Target & $R^2$ & PICP & MPIW \\
\midrule
Energy-Heating & 0.956 & 0.971 & 9.433 \\
Energy-Cooling & 0.922 & 0.964 & 12.524 \\
Concrete & 0.787 & 0.957 & 32.932 \\
\bottomrule
\end{tabular}
\end{table}


WIS and CRPS are evaluated for methods that provide predictive samples or a
parametric predictive distribution.  They are not assigned to the
split-conformal baseline because a conformal interval is not itself a
posterior-predictive distribution.

\FloatBarrier

\section{Rule Selection in the Main Experiment}

\begin{table}[htbp]
\centering
\caption{Mean active-rule ratio and Gibbs effective sample size in the main experiment.}
\label{tab:active}
\scriptsize
\begin{tabular}{lrrr}
\toprule
Target & active-rule ratio & active rules ($R=5$) & ESS/s ($\sigma^2$) \\
\midrule
Energy-Heating & 0.713 & 3.57 & 1563.7 \\
Energy-Cooling & 0.700 & 3.50 & 1599.7 \\
Concrete & 0.800 & 4.00 & 1162.7 \\
\bottomrule
\end{tabular}
\end{table}


The mean active-rule ratio is below one in all three targets, showing that
the rule-level sampler does perform nontrivial selection.  The main-text
comparison nevertheless shows no corresponding predictive improvement over
the dense TSK reference.

\FloatBarrier

\section{Ablation of the Analytical Comparator}

\begin{table}[htbp]
\centering
\caption{Ablation of the analytical model-selection pipeline on Energy-Cooling (10 splits, $R=5$).}
\label{tab:ablation}
\scriptsize
\begin{tabular}{lrrr}
\toprule
Configuration & $R^2$ & PICP & MPIW \\
\midrule
BIC BMA & 0.937 & 0.942 & 9.073 \\
BIC PIP threshold Laplace & 0.934 & 0.938 & 9.239 \\
Gibbs PIP threshold Laplace & 0.933 & 0.939 & 9.390 \\
Full BIC as implemented & 0.931 & 0.936 & 9.537 \\
\bottomrule
\end{tabular}
\end{table}


The four configurations remain close after the corrected membership,
prior, and numerical-stability choices.  In particular, none exhibits the
negative-$R^2$ collapse reported by the earlier implementation.  This
ablation is therefore evidence against treating the earlier failure as a
portable property of BIC, hard thresholding, or Laplace prediction in
isolation.

\FloatBarrier

\section{Coefficient-Level Slab-Variance Sensitivity}

\begin{table}[htbp]
\centering
\caption{Coefficient-level SSVS sensitivity to the slab variance on Energy-Cooling (10 splits).}
\label{tab:tau2}
\scriptsize
\begin{tabular}{lrrrr}
\toprule
$\tau^2$ & $R^2$ & RMSE & PICP & MPIW \\
\midrule
0.1 & 0.894 & 3.093 & 0.969 & 17.739 \\
0.3 & 0.905 & 2.932 & 0.960 & 13.903 \\
1.0 & 0.909 & 2.857 & 0.942 & 11.973 \\
3.0 & 0.917 & 2.730 & 0.944 & 11.157 \\
10.0 & 0.927 & 2.561 & 0.946 & 10.522 \\
100.0 & 0.931 & 2.484 & 0.937 & 9.663 \\
\bottomrule
\end{tabular}
\end{table}


The diffuse end of the grid approaches the dense 10-split reference, while
the tighter settings reduce point accuracy and widen the interval.  The
grid is a sensitivity analysis, not a selection of a new method or a
claim of an optimal prior scale.

\FloatBarrier

\section{Irrelevant-Feature Probe}

\begin{table}[htbp]
\centering
\caption{Irrelevant-feature ablation on Energy-Cooling. Each row appends standard-normal noise features to the eight original inputs.}
\label{tab:noise}
\scriptsize
\begin{tabular}{lrrr}
\toprule
Appended features & TSK-LS & Rule-Gibbs & Coefficient-SSVS \\
\midrule
0 & 0.932 & 0.931 & 0.909 \\
4 & 0.882 & 0.882 & 0.881 \\
12 & 0.878 & 0.878 & 0.879 \\
30 & 0.875 & 0.875 & 0.878 \\
\bottomrule
\end{tabular}
\end{table}


Standard-normal noise features are appended to the eight original
Energy-Cooling features.  The dense, rule-level, and coefficient-level
models remain close across the tested noise levels, so this controlled
probe does not establish a stable advantage for sparsity.

\FloatBarrier

\section{High-Dimensional Probe}

\begin{table}[htbp]
\centering
\caption{High-dimensional probe on Superconductivity (subsample $n=3000$, $d=81$, $R=5$, 30 splits).}
\label{tab:highdim}
\scriptsize
\begin{tabular}{lrrrr}
\toprule
Method & RMSE & $R^2$ & PICP & MPIW \\
\midrule
TSK LS & 14.389 & 0.820 & -- & -- \\
Bayesian TSK & 14.292 & 0.822 & 0.921 & 52.884 \\
SpikeSlab Gibbs & 16.432 & 0.764 & 0.921 & 59.868 \\
SSVS Gibbs ($\tau^2=$1) & 14.655 & 0.813 & 0.934 & 58.375 \\
SSVS Gibbs ($\tau^2=$10) & 14.447 & 0.818 & 0.927 & 54.208 \\
\bottomrule
\end{tabular}
\end{table}


\begin{table}[htbp]
\centering
\caption{Coefficient-level SSVS slab-variance grid on the high-dimensional probe (10 splits).}
\label{tab:highdimtau}
\scriptsize
\begin{tabular}{lrrr}
\toprule
$\tau^2$ & $R^2$ & PICP & MPIW \\
\midrule
0.1 & 0.800 & 0.941 & 61.978 \\
0.3 & 0.808 & 0.940 & 60.654 \\
1.0 & 0.812 & 0.933 & 58.215 \\
3.0 & 0.817 & 0.927 & 56.106 \\
10.0 & 0.819 & 0.923 & 54.269 \\
100.0 & 0.819 & 0.917 & 51.992 \\
\bottomrule
\end{tabular}
\end{table}


The Superconductivity experiment is a 3,000-row subsample with $d=81$ and
$R=5$.  The rule-level Gibbs model is below the dense references in point
accuracy in this probe.  The coefficient-level grid approaches the dense
10-split reference at diffuse slab scales but does not exceed it.
This is a bounded probe, not evidence for a universal high-dimensional
sparsity threshold.

\FloatBarrier

\section{MCMC Diagnostics}

\begin{table}[htbp]
\centering
\caption{Three-chain diagnostics for the rule-level Gibbs sampler on representative splits. Raw coefficient and predictive-quantity diagnostics are reported separately.}
\label{tab:mcmc}
\scriptsize
\begin{tabular}{lrrrrr}
\toprule
Target & raw $\hat R_\beta$ max & predictive $\hat R$ max & $\hat R_{\sigma^2}$ & ESS($\sigma^2$) & draws \\
\midrule
Energy-Cooling & 1.061 & 1.004 & 1.001 & 1639.5 & 2000 \\
Concrete & 1.030 & 1.027 & 1.000 & 1740.3 & 2000 \\
\bottomrule
\end{tabular}
\end{table}


The diagnostic table separates raw coefficient chains from predictive
quantities.  The raw coefficient maxima exceed 1.03 on both representative
targets, whereas the predictive maxima are 1.004 and 1.027.  These values
support reporting finite-chain uncertainty and possible mixing risk rather
than claiming that every parameter chain has converged to a strict
threshold.

\FloatBarrier

\section{Synthetic Enumeration Check}

\begin{table}[htbp]
\centering
\caption{Synthetic enumeration check for the rule-inclusion update and BMA predictive moments.}
\label{tab:synthetic}
\scriptsize
\begin{tabular}{lr}
\toprule
Quantity & Maximum absolute or relative discrepancy \\
\midrule
PIP maximum absolute difference & 0.0261 \\
BMA mean maximum absolute difference & 0.0224 \\
BMA median relative variance difference & 0.0027 \\
\bottomrule
\end{tabular}
\end{table}


The synthetic case has $n=600$, $d=2$, and $R=5$, with active rules
$\{0,2\}$.  Exhaustive enumeration is used only for this small verification
case.  The maximum PIP difference is approximately $0.0261$, the maximum
BMA predictive-mean difference is approximately $0.0224$, and the median
relative predictive-variance difference is approximately $0.0027$.
These values support the arithmetic implementation without implying zero
Monte Carlo error in the benchmark experiments.

\FloatBarrier

\section{Scope and Limitations}

The experiment fixes the antecedent construction, the number of rules, the
prior hyperparameters, and the data splits.  It does not evaluate joint
antecedent--consequent selection, alternative shrinkage priors, measured
building-energy data, or larger high-dimensional samples.  The local
working artifacts document the current rebuild; public archive metadata
must be refreshed by the author before any future submission.

\end{document}
